\chapter{Discussion}\label{ch:conclusion}

A task description's value depends on its precision, its reliability, and the demonstrations already available. In our Gaussian model, a reliable description can replace many demonstrations by constraining the task across directions. Uncertain relevance imposes a finite ESS ceiling before demonstrations are observed. Demonstrations can then supply evidence about relevance or constrain directions left unspecified by a partial description, increasing its ESS.

Meta-trained Transformers approach the Bayes-optimal benchmark, while their response to reliability shifts reflects training. With reasoning enabled, most tested pretrained language models produce description weights closer to the Bayesian reference. For the two models tested with an explicit irrelevant-component prior, agreement with posterior weighting extends from numerical predictions to semantic policy constraints. Incomplete or informal probabilistic specifications produce different behavior, making the stated model assumptions consequential to the comparison.

The analysis offers a controlled interpretation of why descriptions sometimes substitute for many examples \citep{lescao2021,honda2025} and sometimes have limited influence \citep{webson2022,gupta2025,gloaguen2026}. It identifies distinct possible causes: low specificity, uncertain relevance, and a learner's limited use of the available information. Establishing which cause explains a practical prompting result requires evaluating that setting directly.

\section{Scope and Limitations}\label{sec:limitations}

The theory assumes Gaussian linear regression, squared error, and a particular mixture model of irrelevance. Learner-specific ESS measures improvement relative to a learner's own baseline and should be read alongside absolute error. The language-model tasks are small and controlled; the main numerical comparison uses $200$ tasks, and the explicit-prior and semantic experiments cover two models. Some reasoning comparisons also differ in serving, decoding, or readout settings. GPT-6 Sol's results without reasoning depend on its agent configuration. Agreement with a Bayesian reference concerns predictions and does not identify the model's internal computation. Natural descriptions with unstated uncertainty, more structured tasks, and richer or heavy-tailed priors \citep{dalal1983,ohagan2012} would test how far these findings extend.

The wording test with reasoning enabled also covers only two models. The original reliability prompts leave the baseline prior unspecified, so the Bayesian reference for them depends on an assumed prior; we report several choices (Chapter~\ref{ch:llm}). The networks are trained only at $d=5$ and, for the higher-dimension experiment, at $d=10$. Because the numerical language-model task is scalar, it has no random-design dimension term and no partial description, and the language models we test are post-trained chat models. Models without reasoning are evaluated through either the distribution readout or the written readout, so comparisons across these models mix differences between models with differences between readouts.

\section{Future Work}\label{sec:future}

We outline four directions for future work. The first two follow from the limitations and findings above, the third from the descriptions studied here, and the fourth from the networks of Chapter~\ref{ch:networks}. None of them is established in this thesis.

\paragraph{Partial descriptions and the dimension term in language models.}
A precise partial description can become more valuable as demonstrations resolve what it leaves unspecified (Proposition~\ref{prop:partial}), whereas the ESS of a reliable complete description falls toward $c-a_0^{-1}$ (Corollary~\ref{cor:ess}). At $d=16$, $a_0=10$, and $c=100$, ESS increases from $8.1$ without demonstrations to $36.2$ with $100$ demonstrations for $k=8$, and falls from $117$ to $109$ for $k=16$ (Table~\ref{tab:partial}). The networks of Chapter~\ref{ch:networks} capture this increase with $k=3$ of $d=5$ directions described, and their ESS at $c=10$ is larger at $d=10$ than at $d=5$. The numerical language-model task, however, has a single unknown mean. A task with several unknown quantities, of which the description states only some, would test whether a language model's ESS has a dimension term, and whether a partial description becomes more valuable to it as demonstrations accumulate.

\paragraph{Updating reliability.}
With reasoning, GPT-6 Sol and Haiku~5.5 update relevance from the demonstrations once the prompt states the baseline prior: at a stated $50\%$ chance of irrelevance, regressing their normalized description weights on the reference's posterior relevance probability gives slopes of $0.96$ and $0.98$. At a stated $10\%$ chance, however, Haiku~5.5's ESS falls to $15.1$, against $27.5$ for the reference and $25.0$ for GPT-6 Sol. The two models' fitted weights agree within $0.01$ in this condition, so Haiku~5.5's shortfall comes from a minority of replies far from the reference, some of which repeat the demonstration. This rerun covers two models, $c=36$, and $n\leq4$, and its descriptions remain relevant. Three questions remain open. Do the other reasoning models update relevance once the baseline prior is stated? Why do some of Haiku~5.5's replies depart so far from the reference at a small stated chance of irrelevance? And do reasoning models also update relevance under actual mixtures of relevant and irrelevant descriptions, which Table~\ref{tab:mixture} evaluates only without reasoning?

\paragraph{Structural descriptions.}
Apart from the semantic task, the descriptions studied here state where the task lies, through a center $m$ with a stated precision. Many instructions instead state what kind of task it is, for example that only three features matter, or that the answer is monotone in an input. Such a description narrows the hypothesis class rather than pointing at one task. One case of this kind is already covered. A description that confines $w$ to a known $s$-dimensional subspace states the other $k=d-s$ directions exactly, so it is a partial description of those directions in the limit $c\to\infty$. A heuristic calculation suggests that such a description keeps its value as demonstrations accumulate. Under a flat baseline prior, it would lower the Bayes-optimal error from about $d/n$ to about $s/n$, so its ESS would be about $n(d/s-1)$, which grows with the number of demonstrations. Structures that do not fix any direction, such as a sparse task whose relevant features are unknown, or a monotone one, are not covered. The definition of ESS and the network and language-model tests apply to them as well, and the semantic task of Section~\ref{sec:semantic} gives a first test of one in a language model: its broad note names a property the policy uses without stating its required value, and with reasoning, Haiku~5.5 and GPT-6 Sol value it at $1.3$ demonstrations at $n = 2$, against $1.4$ for the reference (Table~\ref{tab:filing}). Moreover, a wrong structural description would leave the Bayesian predictor assuming relevance with a bias that no number of demonstrations removes.

\paragraph{Where a network stores its trust.}
The meta-trained networks inherit their trust from training (Chapter~\ref{ch:networks}). A natural next step is to locate the components of a network that hold this trust, and to test by intervention whether changing them changes the network's implied trust without retraining. Such a study would connect the behavioral measure of this thesis to the mechanisms by which networks learn in context.
